\section{Introduction}
\label{sec:intro}

Commercial and open language-model ecosystems create recurring disputes about
model reuse, derivative services, copied outputs, and suspected distillation.
The hardest cases are not cases where no evidence can exist.  Evidence can be
strong when it lives in a retained protocol transcript, a cryptographic proof,
a provider-side trace, model-state access, or preserved metadata.  The hard
case is ordinary natural final text: a user can copy it, a platform can strip
metadata, and a non-cooperative service need not preserve logs or expose
weights.  This paper studies that gap between where strong evidence can live
and what ordinary final text carries.

We call this the \emph{verification substrate gap}.  A substrate is the object
that carries evidence: metadata, a protocol transcript, a proof object,
provider-side event traces, model state, or the final text itself.  These
substrates support different verification properties.  Some are deterministic
but not portable; some are public but easy to spoof; some are strong only when
the provider cooperates.  Our question is therefore not whether evidence can be
stored somewhere.  It can.  Our question is whether natural LLM final text can
serve as a public, deterministic, portable, non-cooperative verification
substrate.

The experiments reported here were originally motivated by natural-output
evidence injection, but the current evidence points to a different and sharper
conclusion.  Provider-side first-divergence token-event traces can carry
recoverable codeword evidence under reviewed controls.  In the tested
trace-bound diagnostics, Qwen recovers protected codewords in 94 of 96 blocks
and Llama recovers them in 96 of 96 blocks, while raw, task-only, wrong-key,
and wrong-payload controls remain at 0 of 96 for each model.  This rules out a
complete absence of model-side controllability in these diagnostics.

The same evidence does not project into deterministic public final-text
verification.  Across the public final-text predicate artifacts, recovered
codeword blocks remain zero.  Shallow predicates can separate some protected
and raw distributions, but separability is not verification.  The Qwen P2
predicate has AUC 0.554676, protected row TPR 0.275476, raw row FPR 0.193455,
and task-only row FPR 0.195060, with zero recovered codeword blocks.  The Llama
P2 predicate has AUC 0.63128, protected row TPR 0.154992, raw row FPR 0.038798,
and zero recovered codeword blocks.  These predicates observe weak final-text
regularities, not deterministic provenance.

The final-text predicates are also attack surfaces.  Once a public predicate
has a nonzero accept region, source-mismatch rows can be found by rejection
sampling, modified by deterministic rewrite-lite transformations, enriched by
a lightweight surrogate, and directly optimized by guided rewrite/graft.  The
strongest guided attack reaches 100 accepted examples among the top 100 ranked
source-mismatch examples for Qwen raw, Qwen task-only, and Llama raw sources.
These accepts are spoofing evidence against the public final-text predicate
substrate.  They are not protected success, not codeword recovery, and not
public text-only verification.

The paper's central thesis is therefore:
\begin{quote}
Natural LLM outputs exhibit a verification substrate gap: provider-side
first-divergence traces can carry recoverable keyed evidence, but that evidence
does not become a reliable public final-text verifier.  Public final-text
predicates that show shallow separability fail to recover codewords and are
searchable, editable, learnable, and directly optimizable by source-mismatch
examples.
\end{quote}

This thesis deliberately avoids several stronger claims.  We do not claim
public text-only verification success, natural evidence success,
phrase-decoder success, cryptographic provenance, sanitizer robustness,
payload diversity, model-family general verification, or an ownership proof.
The positive evidence in this draft is a trace-bound provider-side diagnostic;
the negative evidence is a public final-text observability and spoofing
diagnostic.

Our contributions are:
\begin{itemize}
\item \textbf{A verification-substrate formulation.}  We define verification
axes that are often conflated: public access, deterministic recovery,
portability, natural-output nativeness, non-cooperation, and spoof resistance.
The current evidence supports a cooperative trace-bundle regime, not all axes
simultaneously.
\item \textbf{A first-divergence diagnostic.}  We formalize first-divergence
token events as a tokenizer-native upper-bound diagnostic for provider-side
event-access settings.  The diagnostic measures model-side controllability; it
is not a deployable public final-text verifier.
\item \textbf{An empirical substrate-gap map.}  We report trace-bound recovery
on Qwen and Llama with clean controls, and contrast it with zero public
final-text codeword recovery.
\item \textbf{A public-predicate attack ladder.}  We show that shallow public
final-text predicates are searchable, editable, learnable, and directly
optimizable by source-mismatch examples, while explicitly not interpreting
those accepted examples as protected success.
\item \textbf{A claim-controlled ownership stress test.}  We map seven dispute
scenarios against nine method families.  The current artifacts support only
cooperative provider-side trace-bundle diagnostics; no current public
final-text-only portable scenario succeeds.
\end{itemize}

The rest of the paper is organized as follows.  Section~\ref{sec:related}
positions the work relative to watermarking, fingerprinting, provenance, and
model-state methods.  Section~\ref{sec:setup} defines the verifier taxonomy
and threat model.  Section~\ref{sec:alignment} introduces first-divergence
events.  Sections~\ref{sec:evidence-map}--\ref{sec:verification} describe the
experimental substrates and public-predicate attacks.  Section~\ref{sec:experiments}
reports the current evidence map, and Section~\ref{sec:limitations} states the
substrate displacement principle and limitations.
